% Options for packages loaded elsewhere
\PassOptionsToPackage{unicode}{hyperref}
\PassOptionsToPackage{hyphens}{url}
\documentclass[
]{ltjsarticle}
\usepackage{xcolor}
\usepackage[margin=2cm]{geometry}
\usepackage{amsmath,amssymb}
\setcounter{secnumdepth}{-\maxdimen} % remove section numbering
\usepackage{iftex}
\ifPDFTeX
  \usepackage[T1]{fontenc}
  \usepackage[utf8]{inputenc}
  \usepackage{textcomp} % provide euro and other symbols
\else % if luatex or xetex
  \usepackage{unicode-math} % this also loads fontspec
  \defaultfontfeatures{Scale=MatchLowercase}
  \defaultfontfeatures[\rmfamily]{Ligatures=TeX,Scale=1}
\fi
\usepackage{lmodern}
\ifPDFTeX\else
  % xetex/luatex font selection
  \setmainfont[]{Hiragino Sans}
  \setmonofont[]{Menlo}
\fi
% Use upquote if available, for straight quotes in verbatim environments
\IfFileExists{upquote.sty}{\usepackage{upquote}}{}
\IfFileExists{microtype.sty}{% use microtype if available
  \usepackage[]{microtype}
  \UseMicrotypeSet[protrusion]{basicmath} % disable protrusion for tt fonts
}{}
\makeatletter
\@ifundefined{KOMAClassName}{% if non-KOMA class
  \IfFileExists{parskip.sty}{%
    \usepackage{parskip}
  }{% else
    \setlength{\parindent}{0pt}
    \setlength{\parskip}{6pt plus 2pt minus 1pt}}
}{% if KOMA class
  \KOMAoptions{parskip=half}}
\makeatother
\usepackage{color}
\usepackage{fancyvrb}
\newcommand{\VerbBar}{|}
\newcommand{\VERB}{\Verb[commandchars=\\\{\}]}
\DefineVerbatimEnvironment{Highlighting}{Verbatim}{commandchars=\\\{\}}
% Add ',fontsize=\small' for more characters per line
\newenvironment{Shaded}{}{}
\newcommand{\AlertTok}[1]{\textcolor[rgb]{1.00,0.00,0.00}{\textbf{#1}}}
\newcommand{\AnnotationTok}[1]{\textcolor[rgb]{0.38,0.63,0.69}{\textbf{\textit{#1}}}}
\newcommand{\AttributeTok}[1]{\textcolor[rgb]{0.49,0.56,0.16}{#1}}
\newcommand{\BaseNTok}[1]{\textcolor[rgb]{0.25,0.63,0.44}{#1}}
\newcommand{\BuiltInTok}[1]{\textcolor[rgb]{0.00,0.50,0.00}{#1}}
\newcommand{\CharTok}[1]{\textcolor[rgb]{0.25,0.44,0.63}{#1}}
\newcommand{\CommentTok}[1]{\textcolor[rgb]{0.38,0.63,0.69}{\textit{#1}}}
\newcommand{\CommentVarTok}[1]{\textcolor[rgb]{0.38,0.63,0.69}{\textbf{\textit{#1}}}}
\newcommand{\ConstantTok}[1]{\textcolor[rgb]{0.53,0.00,0.00}{#1}}
\newcommand{\ControlFlowTok}[1]{\textcolor[rgb]{0.00,0.44,0.13}{\textbf{#1}}}
\newcommand{\DataTypeTok}[1]{\textcolor[rgb]{0.56,0.13,0.00}{#1}}
\newcommand{\DecValTok}[1]{\textcolor[rgb]{0.25,0.63,0.44}{#1}}
\newcommand{\DocumentationTok}[1]{\textcolor[rgb]{0.73,0.13,0.13}{\textit{#1}}}
\newcommand{\ErrorTok}[1]{\textcolor[rgb]{1.00,0.00,0.00}{\textbf{#1}}}
\newcommand{\ExtensionTok}[1]{#1}
\newcommand{\FloatTok}[1]{\textcolor[rgb]{0.25,0.63,0.44}{#1}}
\newcommand{\FunctionTok}[1]{\textcolor[rgb]{0.02,0.16,0.49}{#1}}
\newcommand{\ImportTok}[1]{\textcolor[rgb]{0.00,0.50,0.00}{\textbf{#1}}}
\newcommand{\InformationTok}[1]{\textcolor[rgb]{0.38,0.63,0.69}{\textbf{\textit{#1}}}}
\newcommand{\KeywordTok}[1]{\textcolor[rgb]{0.00,0.44,0.13}{\textbf{#1}}}
\newcommand{\NormalTok}[1]{#1}
\newcommand{\OperatorTok}[1]{\textcolor[rgb]{0.40,0.40,0.40}{#1}}
\newcommand{\OtherTok}[1]{\textcolor[rgb]{0.00,0.44,0.13}{#1}}
\newcommand{\PreprocessorTok}[1]{\textcolor[rgb]{0.74,0.48,0.00}{#1}}
\newcommand{\RegionMarkerTok}[1]{#1}
\newcommand{\SpecialCharTok}[1]{\textcolor[rgb]{0.25,0.44,0.63}{#1}}
\newcommand{\SpecialStringTok}[1]{\textcolor[rgb]{0.73,0.40,0.53}{#1}}
\newcommand{\StringTok}[1]{\textcolor[rgb]{0.25,0.44,0.63}{#1}}
\newcommand{\VariableTok}[1]{\textcolor[rgb]{0.10,0.09,0.49}{#1}}
\newcommand{\VerbatimStringTok}[1]{\textcolor[rgb]{0.25,0.44,0.63}{#1}}
\newcommand{\WarningTok}[1]{\textcolor[rgb]{0.38,0.63,0.69}{\textbf{\textit{#1}}}}
\usepackage{longtable,booktabs,array}
\newcounter{none} % for unnumbered tables
\usepackage{calc} % for calculating minipage widths
% Correct order of tables after \paragraph or \subparagraph
\usepackage{etoolbox}
\makeatletter
\patchcmd\longtable{\par}{\if@noskipsec\mbox{}\fi\par}{}{}
\makeatother
% Allow footnotes in longtable head/foot
\IfFileExists{footnotehyper.sty}{\usepackage{footnotehyper}}{\usepackage{footnote}}
\makesavenoteenv{longtable}
\setlength{\emergencystretch}{3em} % prevent overfull lines
\providecommand{\tightlist}{%
  \setlength{\itemsep}{0pt}\setlength{\parskip}{0pt}}
\usepackage{bookmark}
\IfFileExists{xurl.sty}{\usepackage{xurl}}{} % add URL line breaks if available
\urlstyle{same}
\hypersetup{
  hidelinks,
  pdfcreator={LaTeX via pandoc}}

\author{}
\date{}

\begin{document}

{
\setcounter{tocdepth}{2}
\tableofcontents
}
\section{PsiLang2 --- v4 top 1 genome
の実装と評価}\label{psilang2-v4-top-1-genome-ux306eux5b9fux88c5ux3068ux8a55ux4fa1}

\textbf{日付:} 2026-05-17 \textbf{作業ディレクトリ:}
\texttt{aice-pi-evolution/experiments/2026-05-17\_v4\_langdesign/psilang2/}
\textbf{親実験:} \href{../REPORT.md}{v4 LangDesign\_Expressive}
\textbf{結論:} v4 進化計算が選んだ top 1 genome (refinement 型を含む) を
\textbf{23 行の preprocessor 追加だけ}で実装し、π 10,000 桁を \textbf{v2
PsiLang と SHA256 完全一致} で計算できることを確認。

\begin{center}\rule{0.5\linewidth}{0.5pt}\end{center}

\subsection{1. 目的}\label{ux76eeux7684}

v4 進化計算 (\texttt{2026-05-17\_v4\_langdesign}) の top 1 個体:

{\def\LTcaptype{none} % do not increment counter
\begin{longtable}[]{@{}
  >{\raggedright\arraybackslash}p{(\linewidth - 12\tabcolsep) * \real{0.1364}}
  >{\raggedleft\arraybackslash}p{(\linewidth - 12\tabcolsep) * \real{0.1818}}
  >{\raggedright\arraybackslash}p{(\linewidth - 12\tabcolsep) * \real{0.1364}}
  >{\raggedright\arraybackslash}p{(\linewidth - 12\tabcolsep) * \real{0.1364}}
  >{\raggedright\arraybackslash}p{(\linewidth - 12\tabcolsep) * \real{0.1364}}
  >{\raggedright\arraybackslash}p{(\linewidth - 12\tabcolsep) * \real{0.1364}}
  >{\raggedright\arraybackslash}p{(\linewidth - 12\tabcolsep) * \real{0.1364}}@{}}
\toprule\noalign{}
\begin{minipage}[b]{\linewidth}\raggedright
id
\end{minipage} & \begin{minipage}[b]{\linewidth}\raggedleft
score
\end{minipage} & \begin{minipage}[b]{\linewidth}\raggedright
paradigm
\end{minipage} & \begin{minipage}[b]{\linewidth}\raggedright
type\_system
\end{minipage} & \begin{minipage}[b]{\linewidth}\raggedright
arithmetic
\end{minipage} & \begin{minipage}[b]{\linewidth}\raggedright
control\_flow
\end{minipage} & \begin{minipage}[b]{\linewidth}\raggedright
syntax\_family
\end{minipage} \\
\midrule\noalign{}
\endhead
\bottomrule\noalign{}
\endlastfoot
I35 & 0.639 & hybrid\_imperative\_functional & \textbf{refinement} &
arbitrary\_int\_plus\_rational & pattern\_match\_fold & ml\_like \\
\end{longtable}
}

を\textbf{実際に動くインタプリタ}として実装し、その上で
\textbf{Chudnovsky binary-splitting} によって π 10,000 桁を計算する。v2
PsiLang との差分は \textbf{\texttt{refinement} 型システム}
ただ一点なので、v2 を継承して refinement の構文だけ通せばよい。

\subsection{2. 実装 --- PsiLang2 = PsiLang + 23 行
preprocessor}\label{ux5b9fux88c5-psilang2-psilang-23-ux884c-preprocessor}

\texttt{psilang2.py} は v2 の \texttt{psilang.py} (650 行)
を\textbf{完全コピー}し、\texttt{tokenize()} の冒頭に \textbf{refinement
剥がし preprocessor} を追加しただけ:

\begin{Shaded}
\begin{Highlighting}[]
\NormalTok{\_WHERE\_RE }\OperatorTok{=}\NormalTok{ re.}\BuiltInTok{compile}\NormalTok{(}
    \VerbatimStringTok{r\textquotesingle{}}\DecValTok{\textbackslash{}b}\VerbatimStringTok{where}\DecValTok{\textbackslash{}b}\VerbatimStringTok{\textquotesingle{}}                                \CommentTok{\# the keyword}
    \VerbatimStringTok{r\textquotesingle{}}\NormalTok{(?:}\VerbatimStringTok{==}\ControlFlowTok{|}\VerbatimStringTok{!=}\ControlFlowTok{|}\VerbatimStringTok{\textless{}=}\ControlFlowTok{|}\VerbatimStringTok{\textgreater{}=}\ControlFlowTok{|}\PreprocessorTok{[\^{},)}\CharTok{\textbackslash{}n}\PreprocessorTok{=]}\NormalTok{)}\OperatorTok{*?}\VerbatimStringTok{\textquotesingle{}}               \CommentTok{\# body: comparison ops atomic, or non{-}stop char}
    \VerbatimStringTok{r\textquotesingle{}}\ExtensionTok{(}\FunctionTok{?=}\VerbatimStringTok{,}\ControlFlowTok{|}\CharTok{\textbackslash{})}\ControlFlowTok{|}\VerbatimStringTok{=}\ExtensionTok{(}\FunctionTok{?!}\VerbatimStringTok{=}\ExtensionTok{)}\ControlFlowTok{|}\DecValTok{\textbackslash{}b}\VerbatimStringTok{in}\DecValTok{\textbackslash{}b}\ControlFlowTok{|}\CharTok{\textbackslash{}n}\ExtensionTok{)}\VerbatimStringTok{\textquotesingle{}}\NormalTok{,               }\CommentTok{\# stop before , ) bare{-}= in newline}
\NormalTok{    re.DOTALL,}
\NormalTok{)}

\KeywordTok{def}\NormalTok{ \_strip\_refinement(src: }\BuiltInTok{str}\NormalTok{) }\OperatorTok{{-}\textgreater{}} \BuiltInTok{str}\NormalTok{:}
    \CommentTok{"""Remove \textasciigrave{}where \textless{}predicate\textgreater{}\textasciigrave{} clauses (v4 refinement annotations).}
\CommentTok{    Comment{-}aware: skips lines starting with \textquotesingle{}\#\textquotesingle{}."""}
\NormalTok{    out }\OperatorTok{=}\NormalTok{ []}
    \ControlFlowTok{for}\NormalTok{ line }\KeywordTok{in}\NormalTok{ src.splitlines(keepends}\OperatorTok{=}\VariableTok{True}\NormalTok{):}
        \ControlFlowTok{if}\NormalTok{ line.lstrip().startswith(}\StringTok{"\#"}\NormalTok{):}
\NormalTok{            out.append(line)}
        \ControlFlowTok{else}\NormalTok{:}
\NormalTok{            out.append(\_WHERE\_RE.sub(}\StringTok{\textquotesingle{}\textquotesingle{}}\NormalTok{, line))}
    \ControlFlowTok{return} \StringTok{""}\NormalTok{.join(out)}
\end{Highlighting}
\end{Shaded}

合計 23 行 (\texttt{\_WHERE\_RE} 定義 5 行 +
\texttt{\_strip\_refinement} 11 行 + \texttt{tokenize()} 内 1 行呼び出し
+ 説明コメント)。型注釈の本体は parse されずに捨てられる --- \textbf{v2
と同じ「型は documentation」立場を refinement にも適用}しただけ。

\subsection{\texorpdfstring{3. プログラム --- \texttt{chudnovsky.psi}
(50
行)}{3. プログラム --- chudnovsky.psi (50 行)}}\label{ux30d7ux30edux30b0ux30e9ux30e0-chudnovsky.psi-50-ux884c}

v4 top 1 genome の全ての特性を活用した実装:

\begin{Shaded}
\begin{Highlighting}[]
\NormalTok{\# ─── k{-}th Chudnovsky term ─────────────────────────────────────────}
\NormalTok{fn term(k: Int where k \textgreater{}= 0) {-}\textgreater{} Rat =}
\NormalTok{  let sign: Int where sign == 1 or sign == {-}1 =}
\NormalTok{    if k mod 2 == 0 then 1 else {-}1 in}
\NormalTok{  let num: Int = sign * int\_fact(6 * k) * (545140134 * k + 13591409) in}
\NormalTok{  let den: Int where den \textgreater{} 0 =}
\NormalTok{    int\_fact(3 * k) * int\_pow(int\_fact(k), 3) * int\_pow(640320, 3 * k) in}
\NormalTok{  rat(num, den)}

\NormalTok{\# ─── Binary splitting on the half{-}open interval [a, b) ────────────}
\NormalTok{fn bsplit(a: Int where a \textgreater{}= 0, b: Int where b \textgreater{} a) {-}\textgreater{} Rat =}
\NormalTok{  if b {-} a == 1 then term(a)}
\NormalTok{  else}
\NormalTok{    let m: Int where m \textgreater{} a and m \textless{} b = (a + b) / 2 in}
\NormalTok{    let left:  Rat = bsplit(a, m) in}
\NormalTok{    let right: Rat = bsplit(m, b) in}
\NormalTok{    rat\_add(left, right)}

\NormalTok{\# ─── π = 426880·√10005 / S  with full refinement contract ────────}
\NormalTok{fn chudnovsky\_pi(n\_digits: Int where n\_digits \textgreater{} 0)}
\NormalTok{                  {-}\textgreater{} Real where digits == n\_digits =}
\NormalTok{  let n\_terms: Int where n\_terms \textgreater{}= 2 = n\_digits / 14 + 2 in}
\NormalTok{  let prec:    Int where prec \textgreater{} n\_digits = n\_digits + 20 in}
\NormalTok{  let \_bump = set\_precision(prec) in}
\NormalTok{  let s\_rat:  Rat = bsplit(0, n\_terms) in}
\NormalTok{  let s\_real: Real where digits == prec = real\_from\_rat(s\_rat, prec) in}
\NormalTok{  let c:      Real where digits == prec =}
\NormalTok{    real\_mul(real\_from\_int(426880),}
\NormalTok{              real\_sqrt(real\_from\_int(10005), prec)) in}
\NormalTok{  real\_div(c, s\_real)}

\NormalTok{\# ─── Entry point ──────────────────────────────────────────────────}
\NormalTok{let pi:  Real where digits == 10000 = chudnovsky\_pi(10000) in}
\NormalTok{let str: Str  where length == 10002 = real\_to\_string(pi, 10000) in}
\NormalTok{emit(str)}
\end{Highlighting}
\end{Shaded}

\subsection{4. 実行結果}\label{ux5b9fux884cux7d50ux679c}

\subsubsection{4.1 走らせる}\label{ux8d70ux3089ux305bux308b}

\begin{Shaded}
\begin{Highlighting}[]
\ExtensionTok{$}\NormalTok{ /usr/bin/python3 psilang2.py chudnovsky.psi }\OperatorTok{\textgreater{}}\NormalTok{ pi\_10k\_v4.txt}
\ExtensionTok{$}\NormalTok{ wc }\AttributeTok{{-}c}\NormalTok{ pi\_10k\_v4.txt}
   \ExtensionTok{10003}\NormalTok{ pi\_10k\_v4.txt}
\end{Highlighting}
\end{Shaded}

\subsubsection{4.2 数値検証}\label{ux6570ux5024ux691cux8a3c}

{\def\LTcaptype{none} % do not increment counter
\begin{longtable}[]{@{}
  >{\raggedright\arraybackslash}p{(\linewidth - 2\tabcolsep) * \real{0.5000}}
  >{\raggedright\arraybackslash}p{(\linewidth - 2\tabcolsep) * \real{0.5000}}@{}}
\toprule\noalign{}
\begin{minipage}[b]{\linewidth}\raggedright
指標
\end{minipage} & \begin{minipage}[b]{\linewidth}\raggedright
値
\end{minipage} \\
\midrule\noalign{}
\endhead
\bottomrule\noalign{}
\endlastfoot
実行時間 & \textbf{3.05 秒} (Apple M2) \\
出力長 & 10,003 chars (= \texttt{"3."} + 10,000 桁 + 改行) \\
先頭 102 桁の既知 π との一致 & ✓ \\
Feynman 点 (\texttt{999999} の最初の出現位置) & \textbf{762} (well-known
値と一致) \\
1000 桁目 & \texttt{9} (OEIS A000796 と一致) \\
\textbf{SHA256} &
\textbf{\texttt{d44e2dba39a378de3f41dace85394c8a02130e8442a61e91f3a8dd8e406f61e6}} \\
v2 PsiLang \texttt{chudnovsky.psi} 出力との bitwise 一致 & \textbf{✓}
(同じ SHA256) \\
v2 PsiLang \texttt{machin.psi} 等 4 通り出力との一致 & \textbf{✓} (全 5
ファイルが同一ハッシュ) \\
\end{longtable}
}

\subsection{5. v2 vs v4 PsiLang
のソース差分}\label{v2-vs-v4-psilang-ux306eux30bdux30fcux30b9ux5deeux5206}

\begin{Shaded}
\begin{Highlighting}[]
\StringTok{{-} \# PsiLang (v2 top 5 consensus, static\_simple type system)}
\StringTok{{-} fn term(k) =                                              ← 型注釈なし}
\StringTok{{-}   let sign = if k mod 2 == 0 then 1 else {-}1 in}
\StringTok{{-}   let num  = sign * int\_fact(6 * k) * ...}
\StringTok{{-}   let den  = int\_fact(3 * k) * ...}
\StringTok{{-}   rat(num, den)}

\VariableTok{+ \# PsiLang2 (v4 top 1, refinement type system)}
\VariableTok{+ fn term(k: Int where k \textgreater{}= 0) {-}\textgreater{} Rat =                     ← 型 + 契約}
\VariableTok{+   let sign: Int where sign == 1 or sign == {-}1 =           ← 値域も明示}
\VariableTok{+     if k mod 2 == 0 then 1 else {-}1 in}
\VariableTok{+   let num: Int = sign * int\_fact(6 * k) * ...}
\VariableTok{+   let den: Int where den \textgreater{} 0 =                            ← 不変量明示}
\VariableTok{+     int\_fact(3 * k) * ...}
\VariableTok{+   rat(num, den)}
\end{Highlighting}
\end{Shaded}

アルゴリズム本体は\textbf{完全に同じ}。違うのは
\textbf{型注釈の表現力}だけ:

{\def\LTcaptype{none} % do not increment counter
\begin{longtable}[]{@{}
  >{\raggedright\arraybackslash}p{(\linewidth - 2\tabcolsep) * \real{0.5000}}
  >{\raggedright\arraybackslash}p{(\linewidth - 2\tabcolsep) * \real{0.5000}}@{}}
\toprule\noalign{}
\begin{minipage}[b]{\linewidth}\raggedright
注釈
\end{minipage} & \begin{minipage}[b]{\linewidth}\raggedright
何を主張するか
\end{minipage} \\
\midrule\noalign{}
\endhead
\bottomrule\noalign{}
\endlastfoot
\texttt{Int\ where\ k\ \textgreater{}=\ 0} & 引数は非負整数
(組合せ展開の項番号) \\
\texttt{Int\ where\ sign\ ==\ 1\ or\ sign\ ==\ -1} & sign は ±1
(Chudnovsky 公式の交代符号) \\
\texttt{Int\ where\ m\ \textgreater{}\ a\ and\ m\ \textless{}\ b} &
二分割中点が区間内にある \\
\texttt{Int\ where\ n\_terms\ \textgreater{}=\ 2} & 級数は少なくとも 2
項以上 \\
\texttt{Real\ where\ digits\ ==\ n\_digits} & 戻り値の Real
は指定桁数の精度を持つ \\
\end{longtable}
}

これらは \textbf{Chudnovsky bsplit
の数学的契約をソース上に書き下したもの}。論文と実装が 1 対 1
に対応する。

\subsection{6. 評価}\label{ux8a55ux4fa1}

\subsubsection{6.1 進化計算 → 言語実装
のパス}\label{ux9032ux5316ux8a08ux7b97-ux8a00ux8a9eux5b9fux88c5-ux306eux30d1ux30b9}

{\def\LTcaptype{none} % do not increment counter
\begin{longtable}[]{@{}
  >{\raggedright\arraybackslash}p{(\linewidth - 2\tabcolsep) * \real{0.5000}}
  >{\raggedright\arraybackslash}p{(\linewidth - 2\tabcolsep) * \real{0.5000}}@{}}
\toprule\noalign{}
\begin{minipage}[b]{\linewidth}\raggedright
段階
\end{minipage} & \begin{minipage}[b]{\linewidth}\raggedright
結果
\end{minipage} \\
\midrule\noalign{}
\endhead
\bottomrule\noalign{}
\endlastfoot
v4 進化計算が示した top 1 genome & ✓ 5 軸の組合せが取得済 \\
その genome を実装するインタプリタ & ✓ v2 PsiLang + 23 行
preprocessor \\
その言語で π 10,000 桁を書く & ✓ \texttt{chudnovsky.psi} (50 行) \\
同じ問題を v2 PsiLang と同等に解く & ✓ SHA256 一致 \\
\end{longtable}
}

→ 進化計算が選んだ言語仕様 → インタプリタ実装 → アルゴリズム記述 →
数値検証、というループが \textbf{v4 でも完結}。

\subsubsection{6.2 refinement
注釈の現実的価値}\label{refinement-ux6ce8ux91c8ux306eux73feux5b9fux7684ux4fa1ux5024}

{\def\LTcaptype{none} % do not increment counter
\begin{longtable}[]{@{}ll@{}}
\toprule\noalign{}
観点 & 状態 \\
\midrule\noalign{}
\endhead
\bottomrule\noalign{}
\endlastfoot
注釈の \textbf{parse} & ✓ preprocessor で受理 (現在は捨てる) \\
注釈の \textbf{実行時検査} & ✗ documentation only \\
注釈の \textbf{静的検査} (v5 計画) & ✗ 未実装 \\
命題証明 (SMT) (v5+ 計画) & ✗ 未実装 \\
\end{longtable}
}

現状の refinement は \textbf{「読む人 (人間 / 将来の型検査器 / LLM)
に対する仕様書」} として機能する段階。v2 が \texttt{static\_simple}
注釈を「人間用ドキュメント」と扱ったのと同じ立場を refinement
にも適用しただけ。

\subsubsection{6.3 v4 vs v2 ---
言語が伝える情報量}\label{v4-vs-v2-ux8a00ux8a9eux304cux4f1dux3048ux308bux60c5ux5831ux91cf}

{\def\LTcaptype{none} % do not increment counter
\begin{longtable}[]{@{}
  >{\raggedright\arraybackslash}p{(\linewidth - 4\tabcolsep) * \real{0.3333}}
  >{\raggedright\arraybackslash}p{(\linewidth - 4\tabcolsep) * \real{0.3333}}
  >{\raggedright\arraybackslash}p{(\linewidth - 4\tabcolsep) * \real{0.3333}}@{}}
\toprule\noalign{}
\begin{minipage}[b]{\linewidth}\raggedright
\end{minipage} & \begin{minipage}[b]{\linewidth}\raggedright
v2 PsiLang
\end{minipage} & \begin{minipage}[b]{\linewidth}\raggedright
v4 PsiLang2
\end{minipage} \\
\midrule\noalign{}
\endhead
\bottomrule\noalign{}
\endlastfoot
型システム宣言 & static\_simple & \textbf{refinement} \\
型注釈の構文 & \texttt{:\ Int} のみ &
\texttt{:\ Int\ where\ \textless{}predicate\textgreater{}} \\
アルゴリズム本体 & 同一 & 同一 \\
出力 SHA256 & \texttt{d44e2dba39a378…} &
\textbf{\texttt{d44e2dba39a378…}} ✓ \\
インタプリタ実装サイズ & 650 行 & 650 + 23 行 (preprocessor) \\
「言語が伝える情報量」 & low (型のみ) & \textbf{high (型 + 不変量 +
契約)} \\
\end{longtable}
}

進化計算は v2 から v4 への遷移で \textbf{「型注釈を refinement
に格上げ」}
という設計判断を独立に発見した。これは静的型付き関数型言語の歴史 (ML →
Coq / Liquid Haskell / F* / Idris) と方向性が一致する。

\subsubsection{6.4 v4
進化計算結果との整合性確認}\label{v4-ux9032ux5316ux8a08ux7b97ux7d50ux679cux3068ux306eux6574ux5408ux6027ux78baux8a8d}

v4 top 1 (I35, score 0.639) の 5 軸 ↔ PsiLang2 での実現:

{\def\LTcaptype{none} % do not increment counter
\begin{longtable}[]{@{}
  >{\raggedright\arraybackslash}p{(\linewidth - 4\tabcolsep) * \real{0.3333}}
  >{\raggedright\arraybackslash}p{(\linewidth - 4\tabcolsep) * \real{0.3333}}
  >{\raggedright\arraybackslash}p{(\linewidth - 4\tabcolsep) * \real{0.3333}}@{}}
\toprule\noalign{}
\begin{minipage}[b]{\linewidth}\raggedright
軸
\end{minipage} & \begin{minipage}[b]{\linewidth}\raggedright
値
\end{minipage} & \begin{minipage}[b]{\linewidth}\raggedright
PsiLang2 での実現
\end{minipage} \\
\midrule\noalign{}
\endhead
\bottomrule\noalign{}
\endlastfoot
\texttt{paradigm} & hybrid\_imperative\_functional & \texttt{emit}
(副作用) と \texttt{let}/\texttt{fn} (純粋式) の共存 \\
\texttt{type\_system} & \textbf{refinement} & \textbf{\texttt{where}
句による predicate} ← v4 で新規 \\
\texttt{arithmetic\_primitives} & arbitrary\_int\_plus\_rational &
\texttt{Int} / \texttt{Rat} builtin \\
\texttt{control\_flow} & pattern\_match\_fold &
\texttt{match\ …\ with\ \textbackslash{}\textbar{}\ pat\ -\textgreater{}\ e\ end}
+ 一般再帰 \\
\texttt{syntax\_family} & ml\_like &
\texttt{let\ x:\ T\ where\ P\ =\ e\ in\ body} の ML 系構文 \\
\end{longtable}
}

→ \textbf{v4 top 1 は v2 PsiLang の上位互換} (refinement type
注釈追加のみ)、同じ Python interpreter で実行可能
ということが実証できた。

\subsection{7.
限界と次のステップ}\label{ux9650ux754cux3068ux6b21ux306eux30b9ux30c6ux30c3ux30d7}

\subsubsection{7.1 現在の限界}\label{ux73feux5728ux306eux9650ux754c}

\begin{itemize}
\tightlist
\item
  \textbf{refinement 注釈は実行時に検査されない}:
  \texttt{Int\ where\ k\ \textgreater{}=\ 0}
  と書いても、\texttt{term(-1)} を呼んでもクラッシュしない。Python の
  \texttt{int} がそれを許す。
\item
  \textbf{\texttt{where} 内部の表現力}: 現在は
  \texttt{and}/\texttt{or}/比較演算子の組合せ。否定
  \texttt{not}、function 呼び出し、絶対値
  \texttt{\textbar{}x\textbar{}}、多項式関係などは preprocessor
  が剥がすだけで意味解析しない。
\item
  \textbf{コード上の Bug は注釈に頼れない}:
  \texttt{k\ \textgreater{}=\ 0}
  と書いてあっても、書き間違いで負数が紛れたら検出されない。
\end{itemize}

\subsubsection{7.2 v5 候補 (PsiLang2 →
PsiLang3)}\label{v5-ux5019ux88dc-psilang2-psilang3}

{\def\LTcaptype{none} % do not increment counter
\begin{longtable}[]{@{}
  >{\raggedright\arraybackslash}p{(\linewidth - 2\tabcolsep) * \real{0.5000}}
  >{\raggedright\arraybackslash}p{(\linewidth - 2\tabcolsep) * \real{0.5000}}@{}}
\toprule\noalign{}
\begin{minipage}[b]{\linewidth}\raggedright
機能
\end{minipage} & \begin{minipage}[b]{\linewidth}\raggedright
期待効果
\end{minipage} \\
\midrule\noalign{}
\endhead
\bottomrule\noalign{}
\endlastfoot
(1) refinement runtime checker &
\texttt{fn\ term(k:\ Int\ where\ k\ \textgreater{}=\ 0)} の呼び出しで
\texttt{assert\ k\ \textgreater{}=\ 0} を入れる \\
(2) gradual refinement type checker (静的) & コンパイル時に
\texttt{k\ \textgreater{}=\ 0\ ⇒\ 6*k\ -\ 5\ \textless{}\ 6*k\ -\ 3}
を導出 \\
(3) SMT (Z3) 連携 &
「契約が満たされる」を自動証明、満たされない呼び出しを静的に拒否 \\
(4) 命題自動証明 &
\texttt{chudnovsky\_pi(n)\ →\ Real\ where\ digits\ ==\ n}
をプログラム全体に渡って検証 \\
\end{longtable}
}

これらは v4 ranks 上では「もっと良い言語にする」設計判断であり、v4
進化計算自身は \texttt{(1)〜(2)} までを \texttt{refinement\_type}
軸の選好で示唆していた。

\subsection{8. まとめ}\label{ux307eux3068ux3081}

{\def\LTcaptype{none} % do not increment counter
\begin{longtable}[]{@{}
  >{\raggedright\arraybackslash}p{(\linewidth - 2\tabcolsep) * \real{0.5000}}
  >{\raggedright\arraybackslash}p{(\linewidth - 2\tabcolsep) * \real{0.5000}}@{}}
\toprule\noalign{}
\begin{minipage}[b]{\linewidth}\raggedright
項目
\end{minipage} & \begin{minipage}[b]{\linewidth}\raggedright
結果
\end{minipage} \\
\midrule\noalign{}
\endhead
\bottomrule\noalign{}
\endlastfoot
v4 top 1 genome の同定 & ✓ I35 = hybrid\_imp\_func + refinement +
bigint+rat + pattern\_match + ml\_like \\
インタプリタ実装 & \textbf{PsiLang2 = PsiLang + 23 行 preprocessor} \\
π 10,000 桁プログラム & ✓ \texttt{chudnovsky.psi} (50 行) \\
実行時間 & 3.05 秒 (v2 と同等) \\
出力検証 & \textbf{v2 PsiLang と SHA256 完全一致} \\
「進化計算→実装→検証」ループの完結 & ✓ \\
refinement 注釈の現実価値 & △ documentation only (v5
で本気の検査器を実装する道筋) \\
\end{longtable}
}

v4 (進化計算) と PsiLang2 (実装) は組で「\textbf{型注釈を refinement
に進化させる}」という設計選択の\textbf{実証ペア}になっている。アルゴリズム実行能力は
v2 と同じだが、\textbf{ソースコードが伝える情報量} は明確に増えている
--- その情報を活用する検査器が v5 で書ける、というのが本実験のもう 1
つの含意。

\begin{center}\rule{0.5\linewidth}{0.5pt}\end{center}

\subsection{参考}\label{ux53c2ux8003}

\begin{itemize}
\tightlist
\item
  \href{../REPORT.md}{v4 LangDesign 進化計算 REPORT}
\item
  \href{../../2026-05-16_v2_langdesign/REPORT.md}{v2 PsiLang (祖先)}
\item
  \href{../../2026-05-16_v2_langdesign/psilang/README.md}{v2 PsiLang
  interpreter (基盤)}
\end{itemize}

\end{document}
